% Transposed from accepted baseline 6c16af4; proof is included in full.
\section{Self-contained seam theorem and proof}\label{app:seam}
\subsection{Definitions and statement}

All surrounding circles are externally tangent to a central circle of radius
$R>0$. For surrounding radii $a,b>0$, define

\[
 \theta_R(a,b)=2\arcsin\sqrt{\frac{ab}{(R+a)(R+b)}}\in(0,\pi).
 \tag{A.1}\label{eq:seam1}
\]

If their centers have counterclockwise angular separation $A\in[0,2\pi]$,
the squared distance between them is
$(R+a)^2+(R+b)^2-2(R+a)(R+b)\cos A$. Thus their interiors are disjoint
if and only if

\[
 \theta_R(a,b)\le A\le2\pi-\theta_R(a,b).
 \tag{A.2}\label{eq:seam2}
\]

Indeed, the distance requirement is equivalent to
$\cos A\le1-2ab/((R+a)(R+b))=\cos\theta_R(a,b)$. In particular,
\textbf{both} directed arcs between a pair must have length at least its required
angle. Tangencies, including nonadjacent tangencies, are allowed.

Fix integers $k\ge1$, $n\ge k+2$, and put $N=n-k+1$. The
\textbf{canonical Supnick cycle} $\sigma_{k,n}$ is defined by the following
rank rule, which specifies every entry, without an order ellipsis.
Put $H=\lceil N/2\rceil$. For successive integers $j\ge0$:

\begin{itemize}
\item In a list $A$, append $1+2j$ if it is at most $H$, then append
$N-1-2j$ if it is greater than $H$.
\item In a list $B$, append $2+2j$ if it is at most $H$, then append
$N-2-2j$ if it is greater than $H$.
\end{itemize}

Stop each list when neither entry can be appended. Concatenate $A$, the
reverse of $B$, and the singleton $N$; replace each rank $i$ by
$k+i-1$, and join the last entry to the first. Low ranks partition
$1,\ldots,H$ by parity, high ranks partition $H+1,\ldots,N-1$, and
the appended rank is $N$. Hence this is a cycle through each radius once.
Rotations and reversal represent the same undirected cycle.
Subscripts on cycle entries are read modulo $N$.

Let $E_{k,n}$ be its undirected edge set, and define

\[
 C_{k,n}(R)=\sum_{(a,b)\in E_{k,n}}\theta_R(a,b),\qquad
 C_{k,n}(R_{k,n})=2\pi,\qquad R_{k,n}=R_{\rm chain}(\sigma_{k,n}).
 \tag{A.3}\label{eq:seam3}
\]

Existence and uniqueness of this positive root are proved below. At any
radius write

\[
 \delta_R=\theta_R(n,k)+\theta_R(k,n-1)-\theta_R(n,n-1),\qquad
 \Delta_{k,n}=\delta_{R_{k,n}}.
 \tag{A.4}\label{eq:seam4}
\]

A \textbf{fully feasible placement in this fixed order} consists of angles
$\phi_0<\cdots<\phi_{N-1}<\phi_0+2\pi$, with centers
$(R+s_i)(\cos\phi_i,\sin\phi_i)$, where
$(s_0,\ldots,s_{N-1})=\sigma_{k,n}$, such that \eqref{eq:seam2} holds for every
pair. The \textbf{cumulative-angle placement at the chain root} has $\phi_0=0$
and successive gaps $\theta_{R_{k,n}}(s_i,s_{i+1})$, including the
closing gap. Define $R_{\rm full}(\sigma)$ as the infimum of radii
admitting a fully feasible placement in order $\sigma$.

\textbf{Theorem (fixed-order seam criterion and persistence).} For every pair of
integers $k\ge1,n\ge k+2$:

\begin{enumerate}
\item The positive chain root exists uniquely. The following are equivalent:
$\Delta_{k,n}\ge0$; the cumulative-angle placement at $R_{k,n}$
is fully feasible; some fully feasible placement of this order exists at
$R_{k,n}$. When these hold,
$R_{\rm full}(\sigma_{k,n})=R_{k,n}$.
\item For each fixed $k$, there is a finite integer
$s_k\ge4k+1$ such that $\Delta_{k,n}<0$ for every $n\ge s_k$.
For $k+2\le n<s_k$ the deficit is positive, with at most one
possible exception: $\Delta_{k,s_k-1}=0$. Thus feasibility at the
chain root holds exactly before $s_k$, and failure persists thereafter.
\item For $k=1,2,3$, the complete classifications are:\nopagebreak[4]

\[\begin{array}{c|c|c}k & \Delta_{k,n}>0 & \Delta_{k,n}<0\\1&3\le n\le7&n\ge8\\2&4\le n\le12&n\ge13\\3&5\le n\le16&n\ge17\end{array}\]

In particular $s_1=8,s_2=13,s_3=17$, with no integer equality
case for these three values of $k$.
\end{enumerate}

The proof does not assert the absence of equality for every $k\ge4$,
nor give an explicit formula for their onsets. It establishes the weak
criterion in part 1 even when equality is not excluded.

\subsection{Exact edges, both parities, and growth of the chain root}

Put $L=k+n$. Reading the rank rule gives the following edge lists;
indexed families are listed in increasing order of $i$:

\[
\begin{array}{ll}
 N=2h:\quad &(k,n),\ (k+h-1,k+h),\\
 &(i,L-1-i),\quad k\le i\le k+h-2,\\
 &(i,L+1-i),\quad k+1\le i\le k+h-1;\\[2mm]
 N=2h+1:\quad &(k,n),\\
 &(i,L-1-i),\quad k\le i\le k+h-1,\\
 &(i,L+1-i),\quad k+1\le i\le k+h.
\end{array}
 \tag{A.5}\label{eq:seam5}
\]

Here $h\ge2$ in the even case and $h\ge1$ in the odd case. To
check the formula directly, consecutive low/high ranks inside either list
alternate endpoint sums $N$ and $N+2$. After translation these are
$L-1$ and $L+1$. The closing edge is $(k,n)$. At the join of
the two lists, an even-sized cycle has the additional central edge
$(k+h-1,k+h)$; an odd-sized cycle completes the two indexed families
instead. Their lower endpoints run through exactly the ranges in \eqref{eq:seam5}.
The counts are $2+(h-1)+(h-1)=2h$ and $1+h+h=2h+1$.
The endpoint ranges give distinct edges, so none is omitted or repeated.

For $N=3$ the cycle is $(k,k+1,k+2)$. For every $N\ge4$,
the rank list starts with $(1,N-1)$ and ends with $N$.
Consequently the two neighbors of $k$ are exactly $n-1,n$ in both
parities, including $N=3$.

For fixed positive $a,b$, \eqref{eq:seam1} is continuous and strictly decreasing in
$R$, with limits $\pi$ at zero and zero at infinity. Hence
$C_{k,n}$ decreases continuously from $N\pi>2\pi$ to zero.
This proves the asserted existence and uniqueness.

\textbf{Lemma A.1 (growth without an ordering-optimality premise).} For fixed
$k$, $R_{k,n+1}>R_{k,n}$, and $R_{k,n}\to\infty$.

\emph{Proof.} The kernel strictly increases in either surrounding radius, since
$a/(R+a)$ and $b/(R+b)$ do so separately in \eqref{eq:seam1}.
Compare \eqref{eq:seam5} at $n$ and $n+1$, so $L$ increases by one.

\begin{itemize}
\item If $N=2h+1$, match the seam and each indexed edge $(a,b)$ to
$(a,b+1)$. The new even list contains all these edges and the
additional central edge $(k+h,k+h+1)$.
\item If $N=2h$, match the seam and each indexed edge to $(a,b+1)$,
retaining the old central edge $(k+h-1,k+h)$ unchanged. That retained
edge is the last member of the new lower-sum family. The only remaining
edge in the new odd list is $(k+h,k+h+1)$.
\end{itemize}

Thus for every $R>0$, $C_{k,n+1}(R)>C_{k,n}(R)$. Strict decrease
in $R$ implies strict growth of the roots. No deletion of a vertex and
no optimality over other cycles is used.

All endpoints are at least $k$; at the root,
$2\pi\ge2N\arcsin(k/(R_{k,n}+k))$. Inversion on the increasing
branch, followed by $\sin x\le x$, gives

\[
 R_{k,n}\ge k\bigl(\csc(\pi/N)-1\bigr)
             \ge k(N/\pi-1)\longrightarrow\infty
 \quad (k\text{ fixed}). \tag{A.6}\label{eq:seam6}
\]

This completes the proof. \hfill$\square$

\subsection{The least triangle defect and the fan identity}

\textbf{Lemma A.2 (triangle defect, including equality).} For any three distinct
members $a,b,c$ of $\{k,\ldots,n\}$, at every $R>0$,

\[
 \theta_R(a,b)+\theta_R(b,c)-\theta_R(a,c)\ge\delta_R.
 \tag{A.7}\label{eq:seam7}
\]

Equality holds exactly when $b=k$ and $\{a,c\}=\{n-1,n\}$.

\emph{Proof.} Suppress $R$ from the notation. Differentiating \eqref{eq:seam1} on the
positive domain gives

\[
 \theta_1(a,b)=\frac{\sqrt{Rb/a}}{(R+a)\sqrt{R+a+b}}>0,\qquad
 \theta_{12}(a,b)=\frac{\sqrt R}{2\sqrt{ab}(R+a+b)^{3/2}}>0.
 \tag{A.8}\label{eq:seam8}
\]

Exchange endpoints so that $a<c$. Decreasing the middle argument
$b$ to $k$ gives
$\theta(a,b)+\theta(b,c)-\theta(a,c)\ge H(a,c)$, strictly if
$b>k$, where $H(x,z)=\theta(x,k)+\theta(k,z)-\theta(x,z)$.
Analytic values such as $\theta(k,k)$ are legitimate here; they are
not self-pair constraints. For $x,z\ge k$,

\[
 H_1(x,z)=\theta_1(x,k)-\theta_1(x,z)\le0,\qquad
 H_2(x,z)=\theta_2(k,z)-\theta_2(x,z)\le0.
\]

The inequalities are strict respectively when $z>k$ and $x>k$.
Since $a\le n-1$, $c\le n$, $c>k$, and $n-1>k$,
increasing first $a$ to $n-1$, then $c$ to $n$, proves
$H(a,c)\ge H(n-1,n)=\delta_R$. More explicitly the difference is

\[
 \int_a^{n-1}\int_k^c\theta_{12}(x,y)\,dy\,dx
 +\int_c^n\int_k^{n-1}\theta_{12}(x,y)\,dx\,dy.
 \tag{A.9}\label{eq:seam9}
\]

It vanishes exactly when $a=n-1,c=n$. Together with the strict
middle-argument comparison this proves the equality statement. This argument
covers every ordering of the three distinct radii. \hfill$\square$

For a simple path $P=(v_0,\ldots,v_m)$, define its angular slack as
the sum of its $m$ edge angles minus $\theta_R(v_0,v_m)$.
For $m\ge2$, cancellation gives the exact fan identity

\[
 S_R(P)=\sum_{j=1}^{m-1}
 [\theta_R(v_0,v_j)+\theta_R(v_j,v_{j+1})-\theta_R(v_0,v_{j+1})]
 \ge(m-1)\delta_R. \tag{A.10}\label{eq:seam10}
\]

All triples are distinct, so Lemma A.2 applies term by term. For $m=1$
the slack and the empty sum are both zero. If $\delta_R=0$ and
$m\ge3$, at least two distinct vertices serve as middle arguments;
they cannot both equal $k$. Thus the slack is strictly positive.
For $m=2$, equality in \eqref{eq:seam10} occurs only at the seam through $k$,
up to reversal. No triangle inequality has been assumed: \eqref{eq:seam7} permits a
negative lower bound.

\subsection{Full feasibility, closing gap and small cycles}

Work at $R=R_{k,n}$, and abbreviate $\Delta=\Delta_{k,n}$.
For two distinct positions $i<j$, the two cyclic paths are

\[
 P_+=(s_i,s_{i+1},\ldots,s_j),\qquad
 P_-=(s_j,s_{j+1},\ldots,s_{N-1},s_0,\ldots,s_i).
\]

Their edge counts are $d=j-i$ and $N-d$. Both are simple, even
when one uses the closing edge. If $\Delta\ge0$, applying \eqref{eq:seam10}
separately to both paths shows that each angular length is at least
$\theta_R(s_i,s_j)$. Their sum is exactly $2\pi$ by \eqref{eq:seam3}.
Hence \eqref{eq:seam2} holds for every pair in the cumulative-angle placement. It is
unnecessary to decide which direction is the shorter one.

For adjacent endpoints the one-edge slack is zero. Its complement has
exact slack $2\pi-2\theta_R(s_i,s_{i+1})>0$, since $\theta_R<\pi$.
This checks the upper constraint as well as the lower one, for the closing
pair just as for every other adjacent pair. Formula \eqref{eq:seam2} then proves the
claimed Cartesian non-overlap, and the center radii enforce tangency to the
central circle.

Conversely, suppose some placement of this order is feasible at the chain
root. Unroll its angles and let $g_i>0$ be all $N$ consecutive
gaps, with $g_{N-1}=\phi_0+2\pi-\phi_{N-1}$. Adjacent feasibility
gives $g_i\ge\theta_R(s_i,s_{i+1})$. Therefore

\[
 \sum_i[g_i-\theta_R(s_i,s_{i+1})]=2\pi-C_{k,n}(R)=0.
 \tag{A.11}\label{eq:seam11}
\]

Every summand is nonnegative, so every gap, including the closing gap, equals
its adjacent angle. The two-edge path $(n,k,n-1)$, which exists by
Appendix~A.2, consequently has length
$A=\theta_R(n,k)+\theta_R(k,n-1)$. Feasibility of its endpoints
requires $A\ge\theta_R(n,n-1)$, namely $\Delta\ge0$.
If $\Delta<0$, the complementary arc $2\pi-A$ also exceeds the
upper bound $2\pi-\theta_R(n,n-1)$. There is no freedom to repair
this failure by changing gaps at the same radius and order. This proves
all three equivalences in part 1.

At any feasible radius $r$ for this order, summing adjacent constraints
gives $C_{k,n}(r)\le2\pi$, hence $r\ge R_{k,n}$. Feasibility
at $R_{k,n}$ when $\Delta\ge0$ proves the asserted equality of
the full and chain radii.

For clarity the small cycles are explicit:

\begin{itemize}
\item For $N=3$, $\sigma=(k,k+1,k+2)$. All pairs are adjacent; each
two-edge complement has slack $2\pi-2\theta_R(a,b)>0$.
In particular $\Delta=2\pi-2\theta_R(k+2,k+1)>0$. There is no
nonadjacent pair and no equality case.
\item For $N=4$, $\sigma=(k,k+2,k+1,k+3)$. The pair $(k,k+1)$
has the two paths through $k+2$ and through $k+3$; the pair
$(k+2,k+3)$ has the two paths through $k$ and through $k+1$.
These are precisely the four nonadjacent directed constraints. Lemma A.2
bounds all four defects below by $\Delta$, with equality only on the
path through $k$. In particular the seam complement has two edges,
not three. Adjacent complements have the positive slack already computed.
Appendix~A.5 also gives $\Delta>0$, since $n=k+3\le4k$.
\end{itemize}

For any $N\ge4$, even a hypothetical $\Delta=0$ is sufficient:
only the two-edge seam can have zero slack among multi-edge paths, by \eqref{eq:seam10}
and its equality discussion. When $\Delta>0$, the least nonadjacent
directed-path slack is exactly $\Delta$, attained only by that seam
up to reversal. Longer paths have slack at least $2\Delta$; the other
two-edge paths have strictly larger defects by Lemma A.2.

\subsection{Algebraic threshold and eventual persistence}

We derive the threshold identity rather than invoke a circle-packing
formula. With $x=1/R$, $u=1/a$, $v=1/b$, \eqref{eq:seam1} gives

\[
 \tan\frac{\theta_R(a,b)}2
 =\sqrt{\frac{ab}{R(R+a+b)}}=\frac{x}{\sqrt{x(u+v)+uv}}.
 \tag{A.12}\label{eq:seam12}
\]

Set $q=\sqrt{x(u+v)+uv}>0$ and $w=x+u+v+2q>0$. The positive
square-root identities

\[
 \sqrt{xu+(x+u)w}=x+u+q,\qquad
 \sqrt{xv+(x+v)w}=x+v+q
\]

follow by squaring. Also

\[
 (x+u+q)(x+v+q)-x^2=q(2x+u+v+2q)>0.
\]

Consequently the sum of the half-angles for pairs $(a,1/w)$ and
$(1/w,b)$ has tangent $x/q$. Each half-angle lies in
$(0,\pi/2)$; the positive denominator means their sum also lies
there. It therefore equals $\theta_R(a,b)/2$, with no branch
ambiguity. We have proved

\[
 \theta_R(a,1/w)+\theta_R(1/w,b)=\theta_R(a,b).
 \tag{A.13}\label{eq:seam13}
\]

For the seam endpoints put

\[
 \alpha_n=\frac1n+\frac1{n-1},\quad \beta_n=\frac1{n(n-1)},\quad
 P_n(x)=x+\alpha_n+2\sqrt{\alpha_n x+\beta_n}.
 \tag{A.14}\label{eq:seam14}
\]

This is the Descartes pocket curvature, but \eqref{eq:seam13} proves every identity
needed here independently of that interpretation. Strict increase of the
two angles with the inserted radius gives, for every $R>0$,

\[
 \operatorname{sign}\delta_R
 =\operatorname{sign}(k-1/P_n(1/R))
 =\operatorname{sign}(P_n(1/R)-1/k). \tag{A.15}\label{eq:seam15}
\]

The function $P_n$ is continuous and strictly increasing for $x\ge0$,
and tends to infinity. Its value at zero is
$(1/\sqrt n+1/\sqrt{n-1})^2$, strictly decreasing in $n$.
For $k\ge1$, strict decrease of $t\mapsto1/\sqrt t$ gives
\[
 P_{4k}(0)=\left(\frac1{\sqrt{4k}}+\frac1{\sqrt{4k-1}}\right)^2
 >\left(\frac2{\sqrt{4k}}\right)^2=\frac1k,
\]
\[
 P_{4k+1}(0)=\left(\frac1{\sqrt{4k+1}}+\frac1{\sqrt{4k}}\right)^2
 <\left(\frac2{\sqrt{4k}}\right)^2=\frac1k,
\]
where squaring preserves the strict inequalities because all terms are
positive. Thus $P_{4k}(0)>1/k>P_{4k+1}(0)$, and $\delta_R>0$ for every
$R>0$ when $k+2\le n\le4k$; a unique positive crossing
$P_n(x)=1/k$ exists exactly when $n\ge4k+1$.

On this latter domain set $c=1/k$ and
$q_n=\sqrt{\alpha_n c+\beta_n}$. Squaring the crossing equation
produces $c+\alpha_n\pm2q_n$. The plus root is extraneous because
the unsquared equation requires $c-\alpha_n-x\ge0$, while the plus
root gives $-2(\alpha_n+q_n)<0$. Since the unique positive crossing
exists, it must be the minus root

\[
 \kappa_{k,n}=\frac1k+\frac1n+\frac1{n-1}
 -2\sqrt{\frac{2n+k-1}{kn(n-1)}}>0,\qquad
 T_{k,n}=\frac1{\kappa_{k,n}}. \tag{A.16}\label{eq:seam16}
\]

One can also check its branch directly: $P_n(0)<c$ gives
$q_n^2-\alpha_n^2=\alpha_n(c-\alpha_n)+\beta_n>0$, and
$\alpha_n\kappa_{k,n}+\beta_n=(q_n-\alpha_n)^2$; its positive
square root is $q_n-\alpha_n$. From \eqref{eq:seam15},

\[
 \operatorname{sign}\Delta_{k,n}
 =-\operatorname{sign}(R_{k,n}-T_{k,n})\quad(n\ge4k+1).
 \tag{A.17}\label{eq:seam17}
\]

For fixed $x\ge0$, both $\alpha_n$ and $\beta_n$ strictly
decrease with $n$, so $P_{n+1}(x)<P_n(x)$. Evaluating at the
crossing and using strict increase in $x$ shows
$\kappa_{k,n+1}>\kappa_{k,n}$, hence $T_{k,n+1}<T_{k,n}$.
Equation \eqref{eq:seam16} gives $T_{k,n}\to k$. Lemma A.1 now implies that
$D_{k,n}=R_{k,n}-T_{k,n}$ strictly increases to infinity on
$n\ge4k+1$.

Define $s_k$ as the first integer where $D_{k,n}>0$. It exists;
the deficits are negative from then on. Before it they are positive except
possibly at one integer with $D_{k,n}=0$. Such an equality, if present,
must immediately precede $s_k$ by strict increase. The earlier
no-threshold domain has strictly positive deficits. This proves part 2,
including every quantifier, without assuming monotonicity of the raw
angular deficits.

\subsection{Six exact bridges and the three onset classifications}

By \eqref{eq:seam17} and strict growth of $D_{k,n}$, it suffices to establish

\[
\begin{array}{ll}
 R_{1,7}<6<T_{1,7}, & T_{1,8}<51/10<R_{1,8},\\
 R_{2,12}<17<T_{2,12}, & T_{2,13}<14<R_{2,13},\\
 R_{3,16}<32<T_{3,16}, & T_{3,17}<32<R_{3,17}.
\end{array} \tag{A.18}\label{eq:seam18}
\]

Here is all the rational arithmetic needed to verify these bridges.

\subsubsection{Threshold sides, with signs checked before squaring}

For each row let $r$ be its separator, let
$h=1/k+\alpha_n-1/r>0$, and put
$v=4(\alpha_n/k+\beta_n)>0$. Then
$\kappa_{k,n}-1/r=h-\sqrt v$. The margins in Table~\ref{tab:seam2} are strictly
positive rational numbers:

\begin{table}[htbp]\centering\small
\caption{Exact threshold comparisons; all displayed square margins are positive.}
\label{tab:seam2}
\setlength{\tabcolsep}{5pt}
\begin{tabular}{lrrrll}
\toprule
$(k,n)$ & $r$ & $h$ & $v$ & Positive square margin & Conclusion \\
\midrule
$(1,7)$ & $6$ & $8/7$ & $4/3$ & $v-h^2=4/147$ & $r<T_{1,7}$ \\
$(1,8)$ & $51/10$ & $3061/2856$ & $8/7$ & $h^2-v=47737/8156736$ & $T_{1,8}<r$ \\
$(2,12)$ & $17$ & $1381/2244$ & $25/66$ & $v-h^2=239/5035536$ & $r<T_{2,12}$ \\
$(2,13)$ & $14$ & $643/1092$ & $9/26$ & $h^2-v=673/1192464$ & $T_{2,13}<r$ \\
$(3,16)$ & $32$ & $69/160$ & $17/90$ & $v-h^2=671/230400$ & $r<T_{3,16}$ \\
$(3,17)$ & $32$ & $691/1632$ & $3/17$ & $h^2-v=7465/2663424$ & $T_{3,17}<r$ \\
\bottomrule
\end{tabular}
\end{table}

All six rows lie in the positive-threshold domain, so reciprocation is
legitimate. No sign is inferred by squaring quantities of unknown sign.

\subsubsection{Chain sides, including every closing and central edge}

For each edge $e=(a,b)$ in the \textbf{list order of \eqref{eq:seam5}} put
$z_e^2=ab/((r+a)(r+b))$. In Table~\ref{tab:seam3}, the displayed integer
vector divided by $d$ defines the entire vector $(q_e)$.
An upper row certifies $q_e^2-z_e^2>0$ for every edge; a lower row
certifies $z_e^2-q_e^2>0$. All entries are positive. These are finite
rational comparisons, using \eqref{eq:seam5} and the displayed formula for $z_e^2$.

\begin{table}[htbp]\centering\small
\caption{Rational edge bounds, in the order of \eqref{eq:seam5}.}
\label{tab:seam3}
\setlength{\tabcolsep}{5pt}
\begin{tabularx}{\textwidth}{lrllX}
\toprule
$(k,n)$ & $r$ & Direction & $d$ & Numerators of $(q_e)$, in order \\
\midrule
$(1,7)$ & $6$ & upper & 100 & 28,27,34,37,37,41,43 \\
$(1,8)$ & $51/10$ & lower & 1000 & 316,466,307,390,428,414,462,487 \\
$(2,12)$ & $17$ & upper & 1000 & 209,204,236,257,270,276,250,274,291,301,306 \\
$(2,13)$ & $14$ & lower & 1000 & 245,348,240,278,304,320,330,291,320,340,353,361 \\
$(3,16)$ & $32$ & upper & 100 & 17,23,17,19,20,21,22,22,20,21,22,23,24,24 \\
$(3,17)$ & $32$ & lower & 100 & 17,16,18,20,21,22,22,22,19,21,22,23,24,24,24 \\
\bottomrule
\end{tabularx}
\end{table}

For example, the first edge of the first row is $(1,7)$, giving
$(28/100)^2-1/13=12/8125>0$. An even row includes its central
edge as the second entry. The exact minimum square margin in each row is,
respectively,

\[
 \frac{43}{30000},\quad\frac{4333}{26270000},\quad
 \frac{32}{359375},\quad\frac{3}{32500},\quad
 \frac{23}{70000},\quad\frac{1}{5625}.
 \tag{A.19}\label{eq:seam19}
\]

For upper rows we use two elementary estimates. For $0<z<1/2$,

\[
 \arcsin z<\frac{z}{\sqrt{1-z^2}}<\frac{2z}{\sqrt3}<\frac65z.
 \tag{A.20}\label{eq:seam20}
\]

The first comparison integrates the strictly increasing derivative; the
last follows by squaring positive quantities. For $0<u\le1/3$,

\[
 (1+3u^2/5)^2(1-u^2)-1
 =\frac{u^2(5-21u^2-9u^4)}{25}>0.
\]

Indeed the parenthesis decreases as a function of $u^2$ and equals
$23/9>0$ at $u^2=1/9$. Taking positive square roots and
integrating proves

\[
 \arcsin z<z+z^3/5\quad(0<z\le1/3). \tag{A.21}\label{eq:seam21}
\]

For lower rows simply use $\arcsin z>z$ on $(0,1)$. The standard
strict bounds $3<\pi<22/7$ suffice: the lower bound follows from the
inscribed regular hexagon, and the upper has the elementary witness

\[
 \frac{22}7-\pi=\int_0^1\frac{x^4(1-x)^4}{1+x^2}\,dx>0,
\]

obtained by polynomial division and $\int_0^1dx/(1+x^2)=\pi/4$.
Table~\ref{tab:seam4} gives an exact rational bound for every closure comparison.

\begin{table}[htbp]\centering\small
\caption{Exact chain comparisons for the six endpoint bridges.}
\label{tab:seam4}
\setlength{\tabcolsep}{5pt}
\begin{tabularx}{\textwidth}{lXl}
\toprule
$(k,n)$ & Bound on $C_{k,n}(r)/2=\sum_e\arcsin z_e$ & Strict comparison \\
\midrule
$(1,7)$ & $<(6/5)\sum q_e=741/250$, all $q_e<1/2$ & $<3<\pi$ \\
$(1,8)$ & $>\sum q_e=327/100$ & $>22/7>\pi$ \\
$(2,12)$ & $<\sum(q_e+q_e^3/5)=1457520693/500000000$, all $q_e<1/3$ & $<3<\pi$ \\
$(2,13)$ & $>\sum q_e=373/100$ & $>22/7>\pi$ \\
$(3,16)$ & $<\sum(q_e+q_e^3/5)=14885133/5000000$, all $q_e<1/3$ & $<3<\pi$ \\
$(3,17)$ & $>\sum q_e=63/20$ & $>22/7>\pi$ \\
\bottomrule
\end{tabularx}
\end{table}

Strict decrease of $C_{k,n}$ converts these to the root sides of \eqref{eq:seam18}.
Appendices~A.6.1 and A.6.2 therefore prove all six bridges without numerical
approximations to roots or transcendental functions.

\subsubsection{Completing the integer ranges}

For $k=1,2,3$, let $s$ be respectively $8,13,17$. The
two bridges for that $k$ give $D_{k,s-1}<0<D_{k,s}$. Strict
increase of $D_{k,n}$ on $n\ge4k+1$ supplies the same strict
signs throughout the respective earlier and later ranges. On
$k+2\le n\le4k$, Appendix~A.5 already gives strict positivity of
$\Delta$. Equation \eqref{eq:seam17} and the equivalence in Appendix~A.4 now prove
the entire table in part 3. No integer lies between the two endpoints of
each bridge, so no equality case remains for these three $k$.
